\vspace{-2mm}
\section{Conclusion and Future Work}
\label{sec:conclusion}
This paper introduces FILIP, a simple yet generic framework towards fine-grained vision-language pre-training.
By using a token-wise maximum similarity, our method learns fine-grained representation for patches in the images and words in the sentences.
While it achieves competitive results against several large-scale multi-modal pre-training on various downstream tasks, both its architecture and training procedure can still be optimized to improve its performance. In the future, a
more advanced image encoder as well as a well-designed interaction layer can be used to boost the performance.
Furthermore, we can further add more masked language/image loss to support more generation tasks.
To this end, we hope to extend FILIP as a generic and unified interface for solving a large variety of vision-language tasks.

